\section{Odd arithmetic targets invisible to the primewise and CM realizations}

Write \(G=\widehat{\mathrm{GT}}\), \(H=G_{\mathbf Q}\),
\(M=\ker\lambda\), \(N=\ker(\lambda,\rho_2)\), and
\(K=N\cap H\).  Let \(D\subset N\) be the common kernel of
all rational-prime Kummer cocycles.  For every prime \(r\), let
\(\Phi_r\) be the based standard action on the free pro-\(r\)
fundamental group, with its rational tangential base point.
Let \(T,W\) be the signed cubic and harmonic frames on \(M\).
Define the closed subgroup
\begin{equation}\label{eq:odd-mixed-common-kernel}
 J=D\cap\bigcap_r\ker\Phi_r\cap\ker T\cap\ker W,
 \qquad J_H=J\cap H.
\end{equation}
All intersections and kernels in this formula are taken inside \(N\).

Fix an odd prime \(p\).  Retain the arithmetic target basis
\(\kappa_{ij}\in Y_p=H^1(K,\mathbf F_p)^H\), indexed by
\(i<j\) in \(\{0\}\cup\mathcal S_p\), obtained from
\[
 \zeta_{ij}=\gamma_i\smile\gamma_j-D_{ij}
                      \in H^2(U,\mathbf F_p),
 \qquad U=\widehat{\mathbf Z}^{\times}.
\]
Here \(D_{ij}\) is a finite sum of the cyclic-extension classes.
We use the pushout convention for transgression: an invariant
character is sent to the class of its central pushout extension.
Changing the convention changes every \(\kappa_{ij}\) by the
same sign and has no effect on the assertions of independence.

\begin{theorem}\label{thm:odd-target-survives-all-standard-kernels}
For each pair \(i<j\), there is an arithmetic element
\(w_{ij}\in J_H\) such that, for every \(a<b\),
\begin{equation}\label{eq:odd-mixed-evaluation}
                  \kappa_{ab}(w_{ij})=2\,\delta_{ai}\delta_{bj}.
\end{equation}
Consequently restriction
\[
                         Y_p\longrightarrow H^1(J_H,\mathbf F_p)
\]
is injective.  In particular, if \(\chi\in H^1(N,\mathbf F_p)^G\)
factors through the joint rational Kummer, primewise word, and
two-frame realization of \(N\), then \(\chi|_K=0\).

Thus no nonzero odd arithmetic target character can be obtained
solely from those realizations.  This conclusion does not assert
that a target character fails to extend to full GT.
\end{theorem}

\subsection{Central extensions and alternating evaluation}

We first explain why an odd target character can be evaluated on
the entire cyclotomic kernel \(H^{\mathrm{cyc}}=H\cap M\).
Arithmetic inflation of \(\zeta_{ij}\) is zero.  Represent this
class by a continuous central extension
\[
             1\longrightarrow\mathbf F_p\longrightarrow E_{ij}
                         \longrightarrow U\longrightarrow1.
\]
Vanishing of its arithmetic inflation is equivalent to the
existence of a continuous homomorphism
\(\widetilde\lambda:H\to E_{ij}\) lifting \(\lambda\).
Its restriction to \(H^{\mathrm{cyc}}\) is an \(H\)-invariant
character \(\widetilde\kappa_{ij}\).
Two lifts differ by a character of \(H\), and every such
character factors through \(H^{\rm ab}=U\) by
Kronecker--Weber.  Therefore \(\widetilde\kappa_{ij}\) is
independent of the chosen lift.  Its restriction to \(K\) is
\(\kappa_{ij}\), by the naturality and injectivity of arithmetic
transgression and the isomorphism
\(H^2(U,\mathbf F_p)\simeq H^2(A,\mathbf F_p)\).

If \(c\) represents a central-extension class on an abelian
group, the commutator of lifts of \(u,v\) has central coordinate
\(c(u,v)-c(v,u)\).  The cocycle
\(\gamma_a(u)\gamma_b(v)\) represents
\(\gamma_a\smile\gamma_b\).  Every cyclic-extension class
has zero alternating form, because it is pulled back from an
extension of abelian groups.  Hence, for any \(h,h'\in H\),
\begin{equation}\label{eq:odd-mixed-alternation}
 \widetilde\kappa_{ab}([h,h'])
   =\gamma_a(\lambda(h))\gamma_b(\lambda(h'))
      -\gamma_a(\lambda(h'))\gamma_b(\lambda(h)).
\end{equation}
We use \([h,h']=hh'h^{-1}h'^{-1}\).  The correction
\(D_{ab}\), although essential for arithmetic inflation to
vanish, does not affect this alternating formula.

\subsection{Inertia choices}

Fix decomposition groups and inertia groups \(I_r\subset H\)
for rational primes \(r\).  Inertia at \(r\) acts trivially on
roots of unity of order prime to \(r\), and its \(r\)-adic
cyclotomic image is all of \(\mathbf Z_r^\times\).
The latter assertion follows, for example, from total
ramification of the local cyclotomic tower; its reciprocity
description is \cite[I, Theorem~1.1, p.~20]{MilneCFT}.

Choose \(h_0\in I_p\) whose multiplier has \(p\)-component
in \(1+p\mathbf Z_p\) and satisfies \(\gamma_0(h_0)=1\).
For each \(\ell\in\mathcal S_p\), choose
\(h_\ell\in I_\ell\) with \(\gamma_\ell(h_\ell)=1\).
Every other character among the \(\gamma_i\) is zero on
the chosen element.  In concise notation,
\begin{equation}\label{eq:odd-mixed-inertia-coordinates}
                         \gamma_a(\lambda(h_i))=\delta_{ai}.
\end{equation}
The inertia primes attached to distinct indices are distinct:
index zero uses \(p\), while \(p\notin\mathcal S_p\).

All these elements fix \(\mathbf Q(i,\sqrt{-3})\).
Indeed, every \(\ell\in\mathcal S_p\) is different from
two and three.  For \(p\ne3\), inertia at \(p\) also fixes
this field, which is unramified outside two and three.
For \(p=3\), the imposed condition
\(\lambda_3(h_0)\in1+3\mathbf Z_3\) makes \(h_0\) fix
\(\mu_3\), and inertia at three fixes \(\mathbf Q(i)\).
Thus the assertion holds in this case too.

Put \(z_{ij}=[h_i,h_j]\).  The full multiplier of \(z_{ij}\)
is one.  The tangential pro-\(r\) action of \(H\) on the
thrice-punctured line is unramified outside \(r\):
\begin{equation}\label{eq:odd-mixed-unramified-word}
                              \Phi_r(I_s)=1\qquad(s\ne r).
\end{equation}
The primary statement in \cite[\S8.5, pp.~151--152]{HainMatsumoto2003}
is explicitly an action into \(\operatorname{Aut}\) of the
pro-\(r\) group with tangent vector \(\overrightarrow{01}\);
the subsequent passage to the unipotent completion is separate.
Thus \eqref{eq:odd-mixed-unramified-word} holds for the based
action, not merely its outer or unipotent image.
At least one member of the pair \(h_i,h_j\) has trivial
\(\Phi_r\)-image, for every \(r\).  Therefore
\begin{equation}\label{eq:odd-mixed-z-word}
                              \Phi_r(z_{ij})=1\quad\hbox{for all }r.
\end{equation}

Let \(E_4:y^2=x^3-x\) and \(E_3:y^2=x^3+1\).
Over their CM fields their adelic Tate representations have
abelian images: the Tate modules are rank one over the
completed CM orders, whose unit groups are abelian.
The choices \(h_i,h_j\) fix both CM fields.  Consequently
\begin{equation}\label{eq:odd-mixed-z-CM}
                 \rho_{E_4}(z_{ij})=\rho_{E_3}(z_{ij})=1.
\end{equation}
No good-reduction or local surjectivity assertion at the
coefficient prime is required for this step.

\subsection{Kummer cancellation and detection}

Let \(c\in H\) be complex conjugation, with full multiplier
\(\lambda(c)=-1\).  Define
\begin{equation}\label{eq:odd-mixed-witness}
                         w_{ij}=z_{ij}\,c z_{ij}c^{-1}.
\end{equation}
Each rational Kummer cocycle is additive on \(M\) and obeys
\[
                  \rho_q(aga^{-1})=\lambda(a)\rho_q(g)
                                  \qquad(g\in M).
\]
Since \(z_{ij}\in M\), the two summands in
\eqref{eq:odd-mixed-witness} have opposite Kummer values.
Thus \(\rho_q(w_{ij})=0\) for every rational prime \(q\),
and \(w_{ij}\in D\cap H\).  Normality of the word kernels
and \eqref{eq:odd-mixed-z-word} give
\(\Phi_r(w_{ij})=1\) for every \(r\).

The two Tate representations are homomorphisms, so
\eqref{eq:odd-mixed-z-CM} also holds for \(w_{ij}\).
The established signed arithmetic comparison on
\(H^{\mathrm{cyc}}\) is
\[
 W_g=\rho_{E_4}(g),\qquad
 T_g=\kappa_{3^{1/4}}(g)\rho_{E_3}(g).
\]
The quartic Kummer factor is one on \(D\cap H\).
It follows that \(T(w_{ij})=W(w_{ij})=1\), proving
\(w_{ij}\in J_H\).

Finally, \(\widetilde\kappa_{ab}\) is an \(H\)-invariant
character of \(H^{\mathrm{cyc}}\).  Equations
\eqref{eq:odd-mixed-alternation} and
\eqref{eq:odd-mixed-inertia-coordinates} give
\[
 \kappa_{ab}(w_{ij})
  =2\widetilde\kappa_{ab}(z_{ij})
  =2(\delta_{ai}\delta_{bj}-\delta_{aj}\delta_{bi}).
\]
For ordered pairs \(a<b\) and \(i<j\), the last negative
term is zero.  This proves \eqref{eq:odd-mixed-evaluation}.
Because \(p\) is odd, its diagonal entry is nonzero.
Every element of \(Y_p\) is a finite linear combination
of these basis elements, so their evaluations on the
\(w_{ij}\) prove injectivity of restriction to \(J_H\).

The joint realization of \(N\) in the theorem has kernel
exactly \(J\).  Any character factoring through it is zero
on \(J\), and hence its arithmetic restriction is zero
on \(J_H\).  The proved injectivity makes that arithmetic
restriction zero on all of \(K\), as claimed.

For \(p=3\), this applies in particular to both
\(\kappa_{0,7}\) and \(\kappa_{7,13}\).  Each has an
arithmetic witness with trivial rational Kummer data,
trivial action on every free pro-prime completion, and
trivial full signed cubic and harmonic frames.  A full GT
lift of either class, if one exists, must detect this
common kernel and therefore requires additional information.
